\section{Introduction}

Graph neural networks (GNNs) learn graph representations by repeatedly aggregating information from local neighborhoods \cite{gori2005new,kipf2017gcn}. This message-passing view has become a standard tool for graph classification because it combines node attributes, local connectivity, and graph-level structure in a single differentiable model \cite{gilmer2017neural,wu2020comprehensive}. In molecular and protein-oriented benchmarks, graph labels may depend on both local motifs and global organization, so useful intermediate evidence must survive until graph-level pooling \cite{duvenaud2015convolutional,ying2018hierarchical}.

Increasing depth or adding parallel branches can expand model capacity, but it also changes information transport through the network. Deep GNNs may suffer from over-smoothing, where node representations become too similar across propagation steps \cite{li2018deeper,oono2020simple}. They may also suffer from over-squashing, where long-range information is compressed into limited-size messages \cite{alon2020bottleneck,topping2021understanding}. These issues are especially relevant for graph-level prediction because the final readout depends on the quality of the intermediate representations that reach the pooling stage.

Residual connections are a common response to these difficulties. They improve gradient flow and allow earlier representations to be reused, a principle established in deep convolutional networks and later adapted to graph architectures \cite{he2016deep,bresson2017residual}. Other depth-control strategies, including DropEdge, Jumping Knowledge aggregation, and graph-specific normalization, also aim to preserve useful information across propagation steps \cite{rong2019dropedge,xu2018representation}. However, residual connectivity is still commonly treated as a default layer-wise component rather than as the object of study. Most residual graph classifiers reuse information along depth in the same stream.

Once a model contains multiple branches, depth-wise residual reuse is only one possible route. Residual information may also move across neighboring branches at the same depth or through diagonal branch-depth neighborhoods. Existing graph-classification work has explored stronger message-passing operators, hierarchical pooling, attention, and higher-order neighborhood mixing \cite{abu2019mixhop,bianchi2020graph}. These designs improve the expressive or aggregation capacity of GNNs, but they do not isolate the residual-routing question: if several branch states are available at a given depth, should the model reuse only the previous state of the same branch, exchange states laterally, or combine both directions in a local two-dimensional residual neighborhood?

This paper studies residual reuse as a structured design problem on a branch-by-layer grid. The central question is not whether skip connections should be present, but where residual signals should flow. A grid view exposes three controlled neighborhoods: vertical reuse along depth, horizontal exchange across neighboring branches, and matrix-style reuse that combines same-branch, neighboring-branch, and diagonal residual paths. A sparse/gated matrix variant further tests whether matrix residual traffic should be filtered before injection. This framing follows controlled benchmark practice: the goal is not to add many unrelated architectural components, but to make one design axis explicit enough to evaluate \cite{keriven2020benchmark,morris2020tudataset}.

The study is deliberately controlled. It fixes the graph-classification pipeline within each comparison, evaluates the same architecture family under several canonical message-passing operators, and changes only the residual topology and residual-control mechanism. This separation keeps the empirical comparison focused on information-flow design rather than on unrelated changes in readout structure, task-specific feature engineering, or operator-specific tuning. Recent graph models may achieve stronger absolute performance by changing several components simultaneously; here, those components are intentionally held fixed so that the residual topology itself remains identifiable.

Real graph benchmarks are also limited in the range of structural class granularity that they expose. A residual topology may look adequate on binary graph classification while behaving differently when the label space requires finer structural separation. To test this point without adding task-specific feature engineering, we complement the real-data benchmark with controlled synthetic graph families whose class labels are generated by structural parameters. This synthetic setting is not intended to replace real benchmarks; it is a stress test for whether residual routing remains stable as structural classes become more fine-grained.

The work is organized around three research questions:
\begin{enumerate}
    \item \textbf{Does residual topology affect performance under a shared graph-classification protocol?} We compare plain message passing, vertical residual reuse, horizontal branch-wise exchange, local matrix residual reuse, and sparse/gated matrix reuse while holding the remaining pipeline fixed.
    \item \textbf{How do branch count, residual traffic, and mechanisms explain performance changes?} We test whether additional branches improve accuracy monotonically or only within a useful operating region, and we relate these trends to diversity, similarity, gradients, and residual traffic.
    \item \textbf{When is matrix-style residual reuse more effective on real data and controlled synthetic multi-class structures?} We test whether matrix residual traffic should be passed densely or filtered, whether candidate residual-control settings remain valid after full-fold reruns, and whether matrix-style reuse becomes more stable as structural class granularity increases.
\end{enumerate}

The main contributions are:
\begin{itemize}
    \item \textbf{A branch-by-layer formulation of residual graph classifiers.} Vertical, horizontal, matrix, and sparse/gated matrix residual families are expressed as local neighborhoods on the same grid, making their information-flow differences explicit.
    \item \textbf{A controlled real-data benchmark and synthetic multi-class benchmark.} The experiments keep the classifier, readout, fold protocol, and backbone choices matched, and add a structural stress test spanning five controlled graph families and two to eight classes.
    \item \textbf{A mechanism analysis with conservative rerun checks.} Accuracy is analyzed together with branch diversity, branch cosine similarity, branch CKA, gradient norms, and residual activity. Sensitivity scans are used only to identify candidates, and selected settings are rerun across folds before being treated as empirical evidence.
\end{itemize}
